\subsection{特征值与特征向量的局部微分}

隐函数定理也能研究矩阵变化时的特征方向。设实对称矩阵$\bm S_0\in\mathbb R^{d\times d}$有一个简单特征值$\lambda_0$，对应单位特征向量$\bm v_0$。简单表示这个特征值的代数重数为$1$；对实对称矩阵，这等价于其特征空间是一维的。即使规定向量长度为$1$，$\bm v_0$和$-\bm v_0$仍代表同一方向；为了在局部选定一支，我们要求附近的向量$\bm v$满足$\bm v_0^\top\bm v>0$。

把对称矩阵$\bm S$作为参数，选取其上三角元素作为独立坐标，便可把对称矩阵空间视为有限维欧氏空间。把向量$\bm v$和标量$\lambda$作为未知量。特征方程与单位长度条件合起来写成
\[
 F(\bm S,\bm v,\lambda)=
 \begin{pmatrix}\bm S\bm v-\lambda\bm v\\ (\bm v^\top\bm v-1)/2\end{pmatrix}=\bm0.
\]
在$(\bm S_0,\bm v_0,\lambda_0)$处，关于$(\bm v,\lambda)$的导数作用于$(\bm h,\eta)$，其矩阵为
\[
 \begin{pmatrix}\bm S_0-\lambda_0\bm I&-\bm v_0\\
 \bm v_0^\top&0\end{pmatrix}.
\]
若该矩阵把$(\bm h,\eta)$映为零，先用$\bm v_0^\top$左乘第一行，利用对称性及$\bm S_0\bm v_0=\lambda_0\bm v_0$得到$\eta=0$；第二行给出$\bm v_0^\top\bm h=0$。此时$\bm h$既属于$\lambda_0$的特征空间，又与$\bm v_0$正交。由于特征值简单，只能有$\bm h=\bm0$。因此这个方阵可逆，隐函数定理保证在$\bm S_0$附近存在可微的特征值和单位特征向量分支；上述内积为正的条件固定了局部符号。

对称扰动方向$\bm E$使特征值与特征向量分别产生一阶变化$\dot\lambda$和$\dot{\bm v}$。对特征方程与归一化条件求微分，并在基点取值，得到
\[
 \begin{gathered}
(\bm S_0-\lambda_0\bm I)\dot{\bm v}
 =-(\bm E-\dot\lambda\bm I)\bm v_0\\
\bm v_0^\top\dot{\bm v}=0.
 \end{gathered}
\]
用$\bm v_0^\top$左乘第一式，便得到$\dot\lambda=\bm v_0^\top\bm E\bm v_0$。把它代回后，右侧与$\bm v_0$正交；由于$\lambda_0$简单，$\bm S_0-\lambda_0\bm I$限制在$\bm v_0$的正交补上可逆，因此第二式选出了唯一的$\dot{\bm v}$。

这里的结论依赖所选特征值简单。重特征值使上述关于未知量的导数失去可逆性，因此不能直接用这个构造；这并不意味着每一条特征值分支都不可微。研究重根附近时，可以改为考察相应特征子空间，而不是强行指定其中某一个特征向量。
